\section{All common invariant subspaces}\label{sec:subspaces}
We now show that the scalar multipliers, $S$, and $V$ have no proper nonzero closed common invariant subspace. The first step accounts for every closed subspace invariant under the multipliers, without assuming in advance that it is defined fiber by fiber.

\subsection{Multipliers force a measurable field of projections}
Let $M\subseteq H_0$ be a closed subspace invariant under $S$, $V$,
and every scalar multiplier. Since $M_{\bar g}=M_g^*$ belongs to the
same multiplier family, $M$ reduces each multiplier. Its orthogonal
projection $P$ therefore commutes with all of them. The field lemma
\cite[\FieldImportLocator]{CircleCompanion}, applied with the probability
space $(X,\mu)$ and $K=\ell^2(\Z)$, gives
\[
 (Pf)(x)=P_xf(x),\qquad
 M=\{f\in H_0:f(x)\in\ran P_x\text{ for almost every }x\},
\]
where the $P_x$ are orthogonal projections outside one null set and
$x\mapsto P_xv$ is measurable for each $v\in K$. These are exactly the
projection and measurability conclusions needed below; no covariance
property is imposed on this field.

\subsection{Fiber invariance and measurable tail labels}
For this fixed subspace $M$ and its projection field $P_x$, the global invariance identities $(I-P)SP=0$ and $(I-P)VP=0$ imply
\begin{align}
 (I-P_x)S_xP_x&=0,\label{eq:fiber-S}\\
 (I-P_{\theta x})V_xP_x&=0\label{eq:fiber-V}
\end{align}
almost everywhere. To justify the passage, apply each global identity to the constant field $e_j$. In the second identity evaluate the output at $\theta x$ and use probability preservation. The resulting equalities hold for every $j$ on a common conull set and extend to all of $K$ by boundedness. Intersect this set, the projection set, and $X_*$ over all integer iterates of $\theta$. We can henceforth use all these identities on a single invariant conull set. This set is chosen after fixing $M$; no intersection over all candidate subspaces is required.

By Proposition~\ref{prop:fibers} and \eqref{eq:fiber-S}, there is a label
\[
 s(x)\in\Z\cup\{-\infty,+\infty\},\qquad \ran P_x=\Ktail{s(x)}.
\]
It is measurable: for every $j\in\Z$,
\begin{equation}\label{eq:label-measurable}
 \{x:s(x)\le j\}=\{x:\langle P_xe_j,e_j\rangle=1\}
\end{equation}
on the conull set, and the label can be set to $+\infty$ elsewhere. Equation~\eqref{eq:fiber-V} and Proposition~\ref{prop:V} give two implications:
\begin{align}
 s(x)=-\infty&\ \Longrightarrow\ s(\theta x)=-\infty,\label{eq:full-descent}\\
 s(x)\in\Z&\ \Longrightarrow\ s(\theta x)\le s(x)-1.\label{eq:finite-descent}
\end{align}
The first uses $\overline{V_xK}=K$. The second uses $\overline{V_x\Ktail{h}}=\Ktail{h-1}$ and the inclusion $\Ktail{h-1}\subseteq\ran P_{\theta x}$. When the input fiber is zero, there is no restriction on the output label.

\subsection{Strict descent on a probability space}
The final measure argument is useful independently of the construction.
\begin{lemma}\label{lem:descent}
Let $\theta$ be an invertible ergodic probability-preserving transformation of $(Y,\nu)$. Suppose a measurable function $s:Y\to\Z\cup\{-\infty,+\infty\}$ satisfies, almost everywhere,
\[
 s(x)=-\infty\ \Longrightarrow\ s(\theta x)=-\infty,
 \qquad
 s(x)\in\Z\ \Longrightarrow\ s(\theta x)\le s(x)-1.
\]
Then $s=-\infty$ almost everywhere or $s=+\infty$ almost everywhere.
\end{lemma}
\begin{proof}
All inclusions in this proof are modulo null sets. Put $E=\{s=-\infty\}$. The first implication gives $\theta E\subseteq E$. Probability preservation makes this inclusion an equality modulo null sets, so ergodicity gives $\nu(E)\in\{0,1\}$. The case $\nu(E)=1$ is one conclusion.

Suppose $\nu(E)=0$. For each integer $j$, put $F_j=\{s\le j\}$. Both implications give
\[
 \theta F_j\subseteq F_{j-1}\subseteq F_j.
\]
It follows that $\nu(F_{j-1})=\nu(F_j)$. For any fixed $j$, continuity from above now gives
\[
 \nu(F_j)=\nu\left(\bigcap_{r\ge0}F_{j-r}\right)=\nu(E)=0.
\]
The countable union of the $F_j$ is $\{s<+\infty\}$, proving the other conclusion.
\end{proof}

Applying Lemma~\ref{lem:descent} shows that $P_x=I$ almost everywhere or $P_x=0$ almost everywhere. We have therefore proved that the unital complex algebra generated by $S$, $V$, and all scalar multipliers is transitive.
